\chapter*{前言}\markboth{前言}{}
\addcontentsline{toc}{chapter}{前言}

这本书想回答一个朴素的问题：一个国家的钱、价格、利率、债务和就业，是怎样联系在一起的？我们不从一套现成的模型讲起，而是沿着历史往前走。货币、信用和国家几乎同时出现；此后的五千年里，人们一次次遇到旧理论解释不了的事情，又一次次造出新的理论。跟着这条路走一遍，比背下最后的结论更容易真正学会。

\section*{这本书怎样读}\phantomsection\label{front:howto}

全书有一条主线：\textbf{危机—理论—检验}。每个时代都有一场旧理论解释不了的危机，于是出现新理论；新理论又被下一场危机检验。萨伊定律遇到了大萧条，凯恩斯主义遇到了滞胀，理性预期与新凯恩斯框架遇到了 2008 年，低利率时代的共识又遇到了 2021 年的通胀与高债务。每一章开头交代“当时的世界”，结尾用一节“今天还剩什么”说明这个理论后来被吸收、修正或推翻了哪些部分。

读任何一个宏观论断，我们都分四层来拆：
\begin{enumerate}
\item \textbf{会计恒等式}：永远成立，例如部门余额与政府预算约束。它们约束一切理论，却不告诉我们因果方向。
\item \textbf{行为关系}：消费函数、货币需求、菲利普斯曲线等。它们会随制度与预期改变。
\item \textbf{均衡与预期的闭合方式}：价格怎样调整，人们怎样形成预期。大多数流派之争发生在这一层。
\item \textbf{数据与政策观测量}：我们实际看到的数字，以及它们的测量误差。
\end{enumerate}
很多争论其实是对同一个恒等式的不同因果读法。把四层分开，往往就能看清双方真正分歧在哪里。

\section*{全书的结构}\phantomsection\label{front:structure}

全卷 84 章，分为十一个部分。第一部分先准备语言：国民核算、价格与利率、货币与银行、国际收支。第二、三部分走过古代与近世的各个地区；第四至第十部分按时间讲到 2026 年；第十一部分把全书的工具用于判断，并以对未来的展望结束。

\begin{table}[htbp]\centering\small
\caption{全卷的十一个部分}
\begin{tabularx}{\linewidth}{@{}l Y l l@{}}
\toprule
部分 & 名称 & 时期 & 章 \\\midrule
一 & 宏观的语言 & — & 1—5 \\
二 & 古代世界 & 前 3000 年—600 年 & 6—9 \\
三 & 中世纪与近世 & 600—1750 年 & 10—17 \\
四 & 古典世界 & 1750—1914 年 & 18—25 \\
五 & 战争、萧条与凯恩斯革命 & 1914—1945 年 & 26—34 \\
六 & 战后黄金时代：东方与西方 & 1945—1970 年 & 35—43 \\
七 & 滞胀与预期革命 & 1968—1985 年 & 44—52 \\
八 & 大缓和、新共识与转轨 & 1982—2007 年 & 53—63 \\
九 & 金融危机之后 & 2008—2019 年 & 64—72 \\
十 & 疫情通胀与当下 & 2020—2026 年 & 73—79 \\
十一 & 综合与展望 & — & 80—84 \\
\bottomrule
\end{tabularx}
\end{table}

案例刻意保持多样：除了英美，还有汉代与宋明的中国、罗马、印度洋、马里与印加、德川日本，也有纳粹德国、苏联、改革前后的中国、拉美与东亚的新兴市场。一个理论只有在不同制度下都经得起检验，才值得相信；如果它只在某些制度下成立，知道这些制度条件本身就是收获。

\section*{每一章的结构}\phantomsection\label{front:boxes}

每章以一个导引框开始：核心问题、学完本章应当能做到的事，以及要用到的前面章节。正文依次是当时的世界、理论的回答（先直觉、后推导）、例题、数据检验与反例、今天还剩什么，最后是本章要点和自测题。几种框各有用途：
\begin{itemize}
\item \textcolor{intuit!85!black}{\sffamily\bfseries 直觉}：在推导之前，先用一句话或一幅图抓住要点；
\item \textcolor{exmp}{\sffamily\bfseries 例题}：带完整解答，数字可以由仓库里的脚本复算；
\item \textcolor{cmp}{\sffamily\bfseries 对照}：把两个容易混淆的概念并排放，写出相同、不同与一句话的区分（附录 C 汇总）；
\item \textcolor{casec}{\sffamily\bfseries 案例}：一个具体的时间、地点与事件，每章至少一个来自英美以外；
\item \textcolor{warn}{\sffamily\bfseries 常见误区}：一种流行但不对的说法，以及它错在哪里；
\item \textcolor{deep}{\sffamily\bfseries 深入一步}：研究生水平的推导或技术细节，初读可以跳过。
\end{itemize}

\section*{颜色的含义}\phantomsection\label{front:colours}

正文与插图使用同一套颜色，颜色本身带有含义：
\textcolor{outp}{\rule{0.9em}{0.9em}}~产出与实体经济，
\textcolor{pric}{\rule{0.9em}{0.9em}}~价格与通胀，
\textcolor{mone}{\rule{0.9em}{0.9em}}~货币与利率，
\textcolor{fisc}{\rule{0.9em}{0.9em}}~财政与政府，
\textcolor{extn}{\rule{0.9em}{0.9em}}~对外部门与汇率，
\textcolor{inst}{\rule{0.9em}{0.9em}}~制度与数据。
这套颜色取自对色觉差异友好的 Okabe–Ito 配色。

\section*{数字与数据从哪里来}\phantomsection\label{front:traceability}

正文里的每个数字都可以追溯：要么来自仓库 \texttt{data/snapshots/} 中记录了来源、序列代码与下载日期的数据快照，要么由 \texttt{checks/} 目录下的脚本重新算出并断言。模型用公开的开源工具求解，方法都有出处。前现代与计划经济的数据多为后人的重建估计，正文会写明它们的不确定性；第 81 章专门讨论宏观数据的可靠性。书中引用的外文文献与史料，中文都是本书自己的翻译，原文保存在数据快照里。

对于仍有争议的问题，例如现代货币理论、财政主导、中性利率是否回升、产业政策、休克疗法与渐进改革，本书陈述各方最有力的论证与证据，不替读者下结论。

\section*{分册与写作状态}\phantomsection\label{front:books}

全卷分五个分册编印，章号连续：第一分册（第 1—17 章）宏观的语言与前现代世界；第二分册（第 18—34 章）古典与凯恩斯时代；第三分册（第 35—52 章）战后与滞胀；第四分册（第 53—72 章）大缓和与金融危机；第五分册（第 73—84 章及附录）当下、综合与展望。

这是 2026 年 10 月的骨架版：各章已有导引框，正文将逐章写出。写作计划见仓库中的 \texttt{docs/PLAN.md}，进度见 \texttt{docs/STATUS.md}。
